% ─────────────────────────────────────────────────────────────────────
% DIAGRAM 4 — Crystal of Types → S^11 → SIC Overlap (VAE-Vita Pipeline)
% Author: Lando⊗⊙perator
% ─────────────────────────────────────────────────────────────────────
\documentclass[tikz,12pt,border=12pt]{standalone}
\usepackage{amsmath,amssymb}
\usepackage{fontspec}
\newfontfamily{\shavian}{Trabajo}[Path=/home/mrnob0dy666/.fonts/]
\setmainfont{DejaVu Sans}
\usetikzlibrary{positioning,fit,arrows.meta,calc}

\newcommand{\Dod}{{\shavian\char"10466}}
\newcommand{\Dinfty}{{\shavian\char"1045B}}
\newcommand{\Todot}{{\shavian\char"10478}}
\newcommand{\Tbox}{{\shavian\char"10476}}
\newcommand{\Rdag}{{\shavian\char"1047D}}
\newcommand{\Rcat}{{\shavian\char"10451}}
\newcommand{\Ppmsym}{{\shavian\char"10479}}
\newcommand{\Ppm}{{\shavian\char"1046C}}
\newcommand{\Fhbar}{{\shavian\char"10450}}
\newcommand{\Fell}{{\shavian\char"10471}}
\newcommand{\Kslow}{{\shavian\char"10467}}
\newcommand{\Kfast}{{\shavian\char"1047A}}
\newcommand{\Galeph}{{\shavian\char"10454}}
\newcommand{\Gbeth}{{\shavian\char"10472}}
\newcommand{\Cseq}{{\shavian\char"10460}}
\newcommand{\Cand}{{\shavian\char"1045D}}
\newcommand{\Phic}{{\fontspec{DejaVu Sans}\symbol{8857}}}
\newcommand{\Phisub}{{\shavian\char"10462}}
\newcommand{\Htwo}{{\shavian\char"10456}}
\newcommand{\Hmem}{{\shavian\char"10453}}
\newcommand{\Soneone}{{\shavian\char"10459}}
\newcommand{\Snm}{{\shavian\char"10473}}
\newcommand{\Wz}{{\shavian\char"10474}}
\newcommand{\Wzero}{{\shavian\char"10477}}

\begin{document}
\begin{tikzpicture}[
    block/.style={draw,thick,rounded corners,align=center},
    arrow/.style={-{Latex[length=2.5mm,width=1.8mm]},thick},
]

% Title
\node[font=\Large\bfseries,align=center] at (0,8.2) {Crystal of Types $\to$ $S^{11}$ $\to$ SIC Overlap};
\node[font=\large\itshape,align=center,text=gray!60!black] at (0,7.5) {The VAE-Vita recovers the SIC geometry by learning the Crystal};

% Crystal of Types
\node[block,fill=blue!6,minimum width=16em,minimum height=6em] (crystal) at (-8,4.5)
    {\textbf{Crystal of Types}\\[2pt]
     \small 17,280,000\\[-2pt] \small configurations\\[4pt]
     \footnotesize$\langle$\Dod;\Todot;\Rdag;\Ppmsym;\Fhbar;\Kslow;\\
     \phantom{\footnotesize$\langle$}\Galeph;\Cseq;\Phic;\Htwo;\Soneone;\Wz$\rangle$\\[-2pt]
     \footnotesize$\langle$\Dinfty;\Tbox;\Rcat;\Ppm;\Fell;\Kfast;\\
     \phantom{\footnotesize$\langle$}\Gbeth;\Cand;\Phisub;\Hmem;\Snm;\Wzero$\rangle$\\[-2pt]
     \footnotesize$\vdots$};

% Encoder
\node[block,fill=green!8,minimum width=14em,minimum height=3.6em] (enc) at (-2.5,5.5)
    {\textbf{Encoder $\delta$}\\[2pt]
     \small Crystal ordinal $\to$ vMF $(\kappa,\mu)$ on $S^{11}$};

% Latent space
\node[block,fill=purple!8,minimum width=8em,minimum height=8em,circle] (latent) at (2,5.5)
    {\normalsize$\mathbf{S^{11}}$\\[2pt] \small vMF prior};

% Decoder
\node[block,fill=orange!8,minimum width=14em,minimum height=3.6em] (dec) at (6.5,5.5)
    {\textbf{Decoder $\mu$}\\[2pt]
     \small $S^{11}$ latent $\to$ Crystal logits};

% Reconstruction
\node[block,fill=blue!6,minimum width=12em,minimum height=2.8em] (recon) at (6.5,1.5)
    {\textbf{Reconstructed Crystal}\\[2pt] \small$\mu\circ\delta = \text{id}$};

% SIC property
\node[block,fill=yellow!15,minimum width=16em,minimum height=2.8em] (sicprop) at (2,1.5)
    {\textbf{Mean SIC overlap} $\to \dfrac{1}{d+1} = \dfrac{1}{13}$};

% Arrows
\draw[arrow] (crystal.east) -- (enc.west);
\draw[arrow] (enc.east) -- (latent.west);
\draw[arrow] (latent.east) -- (dec.west);
\draw[arrow] (dec.south) -- (recon.north);
\draw[arrow,bend right=35] (recon.west) .. controls (-2,0) .. (crystal.south);
\draw[arrow,dashed] (latent.south) -- (sicprop.north);
\draw[arrow,dashed] (sicprop.west) -- ++(-1,0) node[left,font=\small] {forced by IC};

% Annotation
\node[align=center,font=\small] at (0,-0.8) {
    VAE-Vita experimental result: SIC overlap $\to 0.0834 \approx 1/13$\\[2pt]
    \textbf{Why?} The Crystal of Types IS the SIC-POVM's outcome space. When a neural network\\[-2pt]
    learns the grammar, the latent space \textit{must} converge to $1/(d+1)$ — forced by\\[-2pt]
    informational completeness. This is the structural SIC-POVM, already solved.
};

\end{tikzpicture}
\end{document}
